\chapter*{Introducci\'on}
\addcontentsline{toc}{chapter}{Introducci\'on}
\markboth{Introducci\'on}{}

El f\'utbol profesional ha atravesado en las \'ultimas d\'ecadas una transformaci\'on
que excede lo deportivo. La digitalizaci\'on masiva de los eventos de juego ha
generado un volumen de datos sin precedentes y ha desplazado el an\'alisis desde el
conteo simple hacia variables que estiman la calidad de las acciones, como los goles
esperados (xG), que valoran las oportunidades creadas antes que el marcador final
\cite{xG,MetricasBase,BigData}. En este entorno, predecir el resultado de un partido
deja de ser un ejercicio de intuici\'on y se vuelve un problema de inferencia
estad\'istica: aunque el desenlace individual conserva una varianza irreducible, el
rendimiento agregado de los equipos exhibe regularidades que pueden estimarse y
proyectarse.

El problema central de esa inferencia es asignar a cada enfrentamiento una
distribuci\'on de probabilidad sobre sus tres resultados posibles, victoria local,
empate o victoria visitante, a partir de una estimaci\'on de la fuerza relativa de
los contendientes. Los organismos rectores del f\'utbol mundial abordan esta tarea
mediante sistemas de clasificaci\'on institucionales, como los rankings de la FIFA
\cite{FIFARanking2024}, la UEFA \cite{UEFARanking2024}, la CONMEBOL
\cite{CONMEBOLRanking2021}, la AFC \cite{AFCRanking2024}, la CONCACAF
\cite{CONCACAFRanking2023} y la IFFHS \cite{IFFHSRanking2021}. Estos sistemas, sin
embargo, operan bajo l\'ogicas predominantemente contables o adaptaciones
pragm\'aticas del modelo Elo \cite{Elo1978}: reducen la fuerza de cada equipo a un
escalar determinista, no cuantifican la incertidumbre asociada a esa estimaci\'on,
no tratan el empate como una categor\'ia con probabilidad propia y no incorporan
informaci\'on contextual del encuentro. Cada una de estas limitaciones deja fuera un
aspecto que la evidencia deportiva sugiere relevante para el resultado.

Frente a ese vac\'io, conviene notar que el rendimiento de un equipo no depende
\'unicamente de su historia de resultados. Factores ex\'ogenos al marcador, de
naturaleza econ\'omica, como el valor de mercado de la plantilla o el balance de
transferencias, y ambiental, como la altitud del estadio o las condiciones
clim\'aticas, han sido se\~nalados por la literatura como posibles moduladores del
desempe\~no \cite{SZYMAN1,Iskandaryan2020,West1996}. Determinar si esa informaci\'on
aporta capacidad predictiva genuina, una vez controlada la fuerza intr\'inseca de los
equipos, es una pregunta abierta y constituye uno de los ejes de este trabajo. De
todo lo anterior se desprende la pregunta de investigaci\'on que orienta el estudio:
\emph{\c{}c\'omo estimar probabil\'isticamente el resultado de un partido de f\'utbol
integrando informaci\'on deportiva, econ\'omica y ambiental, de forma que el sistema
sea actualizable secuencialmente y cuantifique la incertidumbre sobre la fuerza de
cada equipo?}

Para responderla se dise\~na un sistema predictivo organizado en cascada, en el que
cada bloque cubre una limitaci\'on del anterior. Un primer bloque de conteo estima
las tasas de goles esperados mediante modelos de \emph{gradient boosting} con
funci\'on objetivo de Poisson y deriva la probabilidad de empate a trav\'es de la
distribuci\'on de Skellam. Un segundo bloque de rating din\'amico, construido sobre
una implementaci\'on propia del algoritmo Glicko-2 modificado de Shelopugin y
Sirotkin \cite{Glicko2Mod}, actualiza secuencialmente la fuerza latente de cada
equipo, incorpora el empate como categor\'ia estoc\'astica end\'ogena y cuantifica la
incertidumbre de la estimaci\'on. Un tercer bloque, un refinador neuronal basado en
un perceptr\'on multicapa, corrige las probabilidades base integrando las covariables
ex\'ogenas del enfrentamiento. El sistema no pretende sustituir a los modelos
existentes, sino articularlos en una \'unica arquitectura coherente y evaluar, con un
protocolo estad\'istico expl\'icito, qu\'e aporta cada componente al desempe\~no
global. En consecuencia, el objetivo general del trabajo es desarrollar ese sistema
en cascada para estimar probabil\'isticamente el resultado de partidos de f\'utbol a
partir de informaci\'on deportiva, econ\'omica y ambiental, garantizando la
actualizaci\'on secuencial y la cuantificaci\'on de la incertidumbre sobre la fuerza
de los equipos, as\'i como evaluar el aporte real de las covariables ex\'ogenas al
desempe\~no predictivo. Este prop\'osito se concreta en implementar el bloque de
conteo que estima las tasas de goles y deriva la probabilidad de empate, construir la
implementaci\'on propia del Glicko-2 modificado con calibraci\'on de la ventaja de
local\'ia por liga mediante m\'axima verosimilitud, dise\~nar el perceptr\'on
multicapa que refina las probabilidades base con las covariables estandarizadas, y
evaluar el sistema con un protocolo de diez corridas pareadas e inferencia
estad\'istica mediante el test $t$ pareado, analizando el aporte incremental de cada
bloque y el impacto de las variables econ\'omicas y ambientales en la calibraci\'on
global y en la detecci\'on de empates.

El estudio se realiza sobre datos de veintitr\'es competiciones de las principales
confederaciones, comprendiendo las temporadas 2017-18 a 2024-25, organizados en tres
conjuntos de datos que permiten aislar el efecto de las m\'etricas deportivas
detalladas frente al de las variables econ\'omicas y ambientales. La evaluaci\'on
privilegia la calidad de la distribuci\'on de probabilidad reportada, medida
mediante la p\'erdida logar\'itmica, sobre el mero acierto de la clase m\'as
probable, una distinci\'on que resulta central para interpretar correctamente el
desempe\~no del modelo ante la clase empate. En cuanto a su organizaci\'on, el resto
del documento se estructura en cuatro cap\'itulos. El Cap\'itulo~\ref{ch:marco}
desarrolla el marco te\'orico que fundamenta los sistemas de rating, el modelamiento
probabil\'istico del resultado, la integraci\'on de covariables y las m\'etricas de
evaluaci\'on, y precisa la brecha metodol\'ogica que motiva el estudio. El
Cap\'itulo~\ref{ch:metodologia} describe la recolecci\'on y el procesamiento de los
datos, la selecci\'on de covariables y la implementaci\'on, entrenamiento y
calibraci\'on de cada bloque del sistema. El Cap\'itulo~\ref{ch:resultados} presenta
los resultados emp\'iricos obtenidos sobre los tres conjuntos de datos y discute sus
implicaciones. Finalmente, el Cap\'itulo~\ref{ch:conclusiones} re\'une las
conclusiones, expone las limitaciones del estudio y plantea las l\'ineas de
investigaci\'on futuras.